%PLEASE COPY THE FILE AND EDIT THE COPY, DO NOT EDIT THIS FILE!!!

\documentclass[12pt,twoside]{article}
\usepackage[a4paper,width=150mm,top=30mm,bottom=30mm,bindingoffset=10mm]{geometry}
\linespread{1}
\usepackage[utf8]{inputenc} %Standard diacritics in Romance languages (accents, umlauts)
\usepackage{times} %Uses Times New Roman font
\usepackage{wasysym}
\usepackage{tabularx}
\usepackage{float}
\usepackage{diagbox}
\usepackage{graphicx}
\graphicspath{ {./images/} }
\usepackage[style=authoryear-icomp,natbib=true,sortcites=true]{biblatex}% natbib=true so we can use natbib commands with biblatex
% Also: authoryear-icomp is used so that when you cite the same author with different years, you get "according to Herberger (2002, 2004)" rather than "according to Herberger (2002), Herberger (2004)"
\addbibresource{refs.bib}

\usepackage{xpatch}
\usepackage{listings}
\usepackage{color}

\definecolor{dkgreen}{rgb}{0,0.6,0}
\definecolor{gray}{rgb}{0.5,0.5,0.5}
\definecolor{mauve}{rgb}{0.58,0,0.82}

\lstdefinestyle{scala}{frame=tb,
language=Scala,
aboveskip=3mm,
belowskip=3mm,
showstringspaces=false,
columns=flexible,
basicstyle={\small\ttfamily},
numbers=none,
numberstyle=\tiny\color{gray},
keywordstyle=\color{blue},
commentstyle=\color{dkgreen},
stringstyle=\color{mauve},
breaklines=true,
breakatwhitespace=true,
tabsize=3}

\lstdefinestyle{cpp}{frame=tb,
language=C++,
aboveskip=3mm,
belowskip=3mm,
showstringspaces=false,
columns=flexible,
basicstyle={\small\ttfamily},
numbers=none,
numberstyle=\tiny\color{gray},
keywordstyle=\color{blue},
commentstyle=\color{dkgreen},
stringstyle=\color{mauve},
breaklines=true,
breakatwhitespace=true,
tabsize=3}

\lstset{style=scala}


\usepackage[tiny]{titlesec}
\titleformat{\subsection}{}{\thesubsection}{1em}{\itshape}
\titleformat{\subsubsection}{}{\thesubsubsection}{1em}{\itshape}
\titlelabel{\thetitle.\quad}

\usepackage{fancyhdr}
\pagestyle{fancy}
\fancyhead{}
\fancyhead[LE]{\thepage \hspace{3mm} \small }
\fancyhead[RE]{\small Nassouh AlOlabi}
\fancyhead[LO]{\small Complex State Management with Immutability}
\fancyhead[RO]{\small  \normalsize \thepage}
\fancyfoot{}
\setcounter{page}{-3} %This sets the initial page at a number other than 1 (in this case 23).

\fancypagestyle{first}{
	\fancyhead{}
	\fancyhead[L]{\small Higher Institute for Applied Sciences and Technology \\https://hiast.edu.sy\\{}}
	\fancyhead[R]{\small IT Dept. 2022, \\15\\{}}
	\fancyfoot{}}
	%don't worry about the headers, we will compile them

\newcommand{\pref}[1]{(\ref{#1})} % If you use \ref{xx}, the reference in the text appears without parentheses: "as we see in 1, ..." instead of "as we see in (1)...". So we create a new command: instead of \ref, call \pref (Parentheses REFerence) which specifies that any cross references to xx appear in parentheses.

\usepackage{tipa} %for IPA
\usepackage{phonrule} %for phonological rules
\usepackage[nocenter]{qtree} %for trees
\usepackage{gb4e} %for examples and glossing

\usepackage[normalem]{ulem} %STRIKETHROUGH TEXT

\usepackage{authblk,etoolbox}
\renewcommand\Authfont{\Large}
\renewcommand\Affilfont{\normalsize}

\makeatletter
% patch \maketitle so that it doesn't center
\patchcmd{\@maketitle}{center}{flushleft}{}{}
\patchcmd{\@maketitle}{center}{flushleft}{}{}
% patch \maketitle so that the font size for the title is normal
\patchcmd{\@maketitle}{\LARGE}{\normalsize}{}{}

\def\maketitle{{%
		\renewenvironment{tabular}[2][]
		{\begin{flushleft}}
			{\end{flushleft}}
		\AB@maketitle}}
\makeatother

\title{\Huge{Complex State Management with Immutability}}
\author{Nassouh AlOlabi - Dr.Yasser Rahal} \affil{nassouh.alolabi@hiast.edu.sy} 
\setlength{\affilsep}{1pt}
\date{}

%IF THERE ARE TWO AUTHORS:
%\title{\Huge{Patterns of syntactic microvariation: the case of European Portuguese}}
%\author{Noam Chomsky} \affil{Massachusetts Institute of Technology\\noam.chomsky@mit.edu} 
%\author{Howard Lasnik} \affil{University of Maryland\\howard.lasnik@umd.edu}
%\setlength{\affilsep}{1pt}
%\date{}

\begin{document}

\pagenumbering{roman}

\begin{figure}
    \centerline{\includegraphics[scale=.25]{logo.png}}
\end{figure}
\maketitle

\thispagestyle{first}

\vspace{0.5cm}

\hfill Received: 22-02-2022 

\hfill Accepted: 25-02-2022

\hfill Published: 05-03-2022

\vspace{1cm}

% \noindent \textbf{How to cite} Leave blank

\vspace{1.5cm}

\noindent \textbf{Abstract}
\begin{center}
 	\line(1,0){430}
\end{center}
\vspace{-0,3cm}
\noindent Recent growing demand for fault-tolerant, scalable, distributed systems has made some mainstream software architectures and patterns obsolete or rather harder to come by, and thus came the rise of stateless and functional solutions based on data immutability which has already been the cornerstone of "Big Data"

\vspace{5mm}

\noindent \textbf{Keywords:} immutable data structures, functional programming, algorithm design

\vspace{4mm}
\begin{center}
	\line(1,0){430}
\end{center}

%\begin{center}
%	\textbf{Table of Contents}
%\end{center}

%\begin{large}
%\begin{center}
	%\begin{tabular}{c c}
%		1. Section 1 & 4. Section 4\\
%		2. Section 2 & 5. Section 5\\
%		3. Section 3 & 6. Section 6
%	\end{tabular}
%\end{center}
%\end{large}

\newpage
\tableofcontents
\newpage

\pagenumbering{arabic}
\setcounter{page}{1}

\section{Introduction}
% immutability an empircal study in scala
Unintended mutations of a program’s state are one reason for inconsistent behavior and bugs of the
program. These mutations might have been introduced by side-effects of functions that developers
were unaware of during the program’s implementation. Some programming languages do, for example,
allow arguments of a function to be mutated. The fact that a third-party function can mutate the
state of its argument can go undetected. The problem of rogue and complicated state mutations can
become difficult to handle when states are shared among objects. One way to avoid undesired mutations
is to use immutable data instead of mutable data. Immutable data cannot be mutated once created
and instead of mutating shared data in memory, data would have to be re-created to include the modifications
needed. Programmers can then safely share data without the possibility of having it mutate to
something else, which is crucial to avoid, for example, race conditions in concurrent programs.

\section{State Management}
% https://medium.com/super-declarative/understanding-state-management-and-why-you-never-will-dd84b624d0e
The term state management has gain traction over the last decade or so with the emergence of wide spread IT solutions along with it the ease of which you can make one of your own from whatever background you come from. And so depending on the underlying scenario the term could stand for:
Data persistence management,
Information flow,
Programming paradigms,
I/O (especially networking and caching), %% TODO: latency table
Application architecture,
Presentation behavior and 
UI templating. So the term state management has been a catch-all term from data modeling to parallel computing. 

\subsection{Data Persistence}

%%-----------https://www.researchgate.net/publication/221596019_Analyzing_persistent_state_interactions_to_improve_state_management ---------
A primary challenge to building reliable and secure computer systems is managing the persistent state of the system: all the executable files, configuration settings and other data that govern how a system functions. The difficulty comes from the sheer volume of this persistent state, the frequency of changes to it, and the variety of workloads and requirements that require customization of persistent state. The cost of not managing a system's persistent state effectively is high: configuration errors are the leading cause of downtime at Internet services, troubleshooting configuration problems is a leading component of total cost of ownership in corporate environments, and malware—effectively, unwanted persistent state—is a serious privacy and security concern on personal computers.
%In this paper, we analyze how computer systems dynamically interact with files and configuration settings in an attempt to gain insights into the problem of persistent state management. We analyze over 3648 machine days of these persistent state interactions, collected over an 8 month period from 193 machines. These machines are under real workloads and include Internet servers, corporate desktops, and home machines. We characterize the scope and magnitude of the persistent state management problem today, measuring not only the gross characteristics of persistent state, but also analyzing how it is used by applications, and when administrators and users modify it. We find that monitoring persistent state interactions provides important visibility and show how it can be used as a foundation for building better persistent state management tools.
%%----------------------------------------------------

\subsection{Information Flow}
%%https://thedaylightstudio.com/blog/2018/03/14/what-is-state-in-web-application-development
Application State (a.k.a. Program State) represents the totality of everything necessary to keep your application running.  When we refer to application state we are normally referring to the state of the program as it exists in the contents of its memory.  But what does that mean practically?  How am I to understand that?  It helps to think in extremes.  What happens to information and functionality core to your application if a server goes down and restarts?  You lose whatever was residing in memory.

%This is one of the reasons why in web development we often use stateless resource controllers which disseminate the information necessary to the running of your application in a way that does not rely heavily on holding data for retrieval on the server’s memory.
%% -------------------------------

%% https://dojotoolkit.org/documentation/tutorials/1.6/data_modeling/
 %The Model-Viewer-Controller (MVC) is a dominant paradigm for application development. The MVC approach separates key common concerns for organized, manageable application code. Dojo is heavily based on MVC principles, and provides powerful helpers for MVC-structured applications. The foundation of a well-designed MVC application is a solid data model. Here we will see how we can leverage Dojo object stores and Stateful objects to create a robust model that can be used in the view and controller code. The model is the M in MVC. The data model represents the core information that your application is being used to access and manipulate. The model is the center of your application
%%--------------------------------------



%%%%https://www.researchgate.net/publication/2716493_The_Role_of_Distributed_State
%Distributed state offers the potential for improving the performance, coherency, and reliability of distributed systems. Unfortunately, distributed state also introduces consistency problems, crash sensitivity, time and space overheads, and complexity; these problems make it difficult to achieve the potential benefits. This paper describes the advantages and disadvantages of distributed state, and presents the NFS and Sprite file systems as examples of different tradeoffs. It does not appear possible to achieve all the advantages of distributed state and also avoid all the problems; rather, system designers must make compromises based on the needs of their individual environments. The Role of Distributed State February 19, 1990 1. Introduction Webster's New World Dictionary defines state as "a set of circumstances or attributes characterizing a person or thing at a given time" [4]. State plays a fundamental role in all computer systems. One way of characterizing computation is as a ...
%----------------

Anti-Patterns or pitfalls includes:

% http://www.padsweb.rwth-aachen.de/wvdaalst/publications/p514.pdf
\subsubsection{Missing data}
Its' the situation where some data element needs to be accessed, i.e. read or destroyed, but either it has never been created or it has been deleted without having been created again.
\subsubsection{Inconsistent Data} 
Data is inconsistent if a task is using this data while some other task (or another instance of the same task) is writing to this data or is destroying it in parallel.
%-----------------------------
\subsection{Programming Paradigms}
% https://www.geeksforgeeks.org/introduction-of-programming-paradigms/
Programming paradigm is an approach to solve problem using some programming language or also we can say it is a method to solve a problem using tools and techniques that are available to us following some approach.
%----------------------------------------
And so most paradigms have their own opinion on data manipulation for example: Procedural \& Object Oriented paradigms allows more access to data modifying (destructive updates), deleting. Others like Logic \& Functional doesn't allow them.

\subsection{I/O (especially networking and caching)}
%https://www.solarwinds.com/-/media/solarwinds/swdcv2/licensed-products/storage-resource-monitor/resources/whitepapers/top_4_causes_of_storage_io_bottlenecks.ashx
Applications that are I/O heavy often cause bottlenecks. As well as often causing the problem, I/O intensive applications are often more sensitive to a storage latency issue. When you have a large user base trying to access these applications, slowdowns tend to take place. Increased response time in storage I/O causes bottlenecks. When there is a queue in the storage I/O, you would generally see an increase in latency.
Network and Disk latency are often times the biggest source of headache if poorly managed, which is often times the case since they're usually treated as if they were function calls which they're clearly not.
% http://norvig.com/21-days.html#answers
\begin{center}
    \begin{table}[H]
        \begin{center}
        \begin{tabularx} {0.8\textwidth}{ 
            | >{\raggedright\arraybackslash}X 
        | >{\raggedleft\arraybackslash}X | }
        \hline
        execute typical instruction &	1/1,000,000,000 sec = 1 nanosec\\
\hline
    fetch from L1 cache memory &	0.5 nanosec\\
\hline
branch misprediction &	5 nanosec\\
\hline
    fetch from L2 cache memory &	7 nanosec\\
\hline
    Mutex lock/unlock &	25 nanosec\\
\hline
    fetch from main memory &	100 nanosec\\
\hline
    send 2K bytes over 1Gbps network &	20,000 nanosec\\
    \hline
    read 1MB sequentially from memory &	250,000 nanosec\\
\hline
    fetch from new disk location (seek) &	8,000,000 nanosec\\
\hline
    read 1MB sequentially from disk &	20,000,000 nanosec\\
\hline
    send packet US to Europe and back &	150 milliseconds = 150,000,000 nanosec\\
    \hline      
    % \float
\end{tabularx}
    
\end{center}
\caption{ \label{tab:1} Latency numbers every programmer should know}
\end{table}
\end{center}
As you can see disk fetches can take 100x more than main memory fetches. And when dealing with distributed systems were network call are an integral part of the system treating network call which 100000x more than main memory calls you end up with catastrophically unpredictable behavior



%-------------------------------------------------------


%this format problems tech solutions
\section{Immutability}


% src
% Immutability an Empirical Study in Scala.pdf
When something is immutable, we say that it cannot be changed or that it is unchangeable. The definition of an immutable object in an object-oriented programming language is an object that has a state that cannot be mutated once instantiated (created). Moreover, an immutable class is a class whose instances cannot be mutated, meaning that there are no methods in the class that can mutate an instance of it. Classes that are not immutable can, however, have instances of it that are immutable. 

\subsection{Forms of Immutability}
There are different varieties of immutability used in practice today, which includes:
\begin{itemize}
    \item Object immutability
    \begin{description}
        \item[] An immutable object is an object that cannot be modified (mutated), i.e., its state cannot be mutated
    \end{description}
    \item Class immutability 
    \begin{description}
        \item[] When every instantiated object of a class is immutable, then we say that the class itself is immutable.
    \end{description}
    \item Deep and shallow immutability (transitivity)\begin{description}
        \item[] Immutability can be deep or shallow, i.e., transitive or non-transitive:
        \item[]Deep (transitive): immutability means that all objects referred to by an immutable object must also be immutable.
        \item[]Shallow (non-transitive) immutability has a more relaxed constraint, and the immutable object may refer to objects that are mutable, but the fields of the immutable object itself cannot be mutated.
    \end{description}    
    \item Reference immutability (read-only references)
     \begin{itemize}
    \item[] Languages with the support of reference immutability have the notion of read-only references. A readonly reference cannot mutate the object it is referring to. When all references to an object are readonly, then nothing can mutate the object, and the object is immutable. There are, however, often no guarantee that a referred to object is immutable because the object may still be mutated by some other reference that is not read-only, unless there is some analysis that ensures that there only exist readonly references to that object. Reference immutability thus often ensure shallow (non-transitive) immutability and “reference immutability” does not mean that a reference itself is immutable.
\end{itemize}


\item Non-assignability
     \begin{itemize}
    \item[] Non-assignability is a form of immutability and property of a variable that makes sure that it cannot be reassigned. Since fields of objects are variables and if no variable can be reassigned after its initial assignment, the object is effectively immutable if what the fields are referring to is immutable. This make non-assignability give shallow (non-transitive) immutability too. Assignment of an object’s field mutates the object, but it does not mutate what was previously on the field or what was assigned to the field.
\end{itemize}


\item Concrete and abstract immutability
     \begin{itemize}
    \item[] Concrete immutability does not allow any change to the object’s state. Abstract immutability, however, allows an immutable object to change its internal state but not the object’s “abstract” value. The object is still immutable from the perspective of the object’s observer that can only see the abstract value, but the object may mutate internally. This can, for example, be useful to speed up certain operations, lazy initialization and buffering.
\end{itemize}


\end{itemize}

\subsection{Properties of Immutability}
% write some problems
The concept of immutability is important and
utilizing immutable data is considered to ease software development and reasoning about programs in
numerous ways, for example:

\subsubsection{Predictability} It is harder to understand and reason about programs that have shared mutable
states with unclear interactions. Tracking mutations and maintaining the correct state of a program
can be difficult. Using immutable data naturally, avoids state changes and forces the programmer
to let data flow and be utilized in a different way throughout the program, making the state of the
program more predictable. For example, calling the same function twice would yield the same result,
and the outcome is predictable. Without the immutability guarantee, the second call could
yield another result because of an underlying state mutation making it less predictable.

\subsubsection{Testability} Because immutable data can only be changed once during construction, they are inherently
simpler and easier to unit test. One may reason that by restricting the number of possible mutations
in a program; the number of potential errors of the program is also reduced. Testing is essentially
to validate that mutations in the program occur correctly and thus having more mutations would require more testing. By restricting the number of mutations in a program, the program has
fewer reasons for errors to occur and there are fewer state transitions to test.

\subsubsection{Concurrency} Immutable data are thread-safe, as data cannot mutate, there is no danger in having
multiple threads access the same data at the same time and have synchronization issues.

\subsubsection{Modularity} Without depending on a local or global state, immutable types and data may be reused
in different contexts more easily.

\section{Methods}

\subsection{Shadowing}

Shadowing is a technique that could represent a changing variable. for instance, an accumulator. the technique is simple and possible in almost every programming language out there.
at it's simplest form shadowing looks like this:
\begin{lstlisting}
    // Scala  
    object Main {
        def main(args: Array[String]){
            val i = 1;
            {
                val i = 2;
                {
                    val i = 3;
                }
            }
        }  
    }    
\end{lstlisting}
But that's not really useful and even more confusing and very error prone and I'd agree shadowing in of it's self isn't really useful but it's really at the core of any recursive solution since every function come with it's own block and 
\begin{lstlisting} 
    def factorial (n: BigDecimal):BigDecimal = {
        if (n <= 1)  1
        else n * factorial(n-1)
    }
\end{lstlisting}
Notice here n value range over \{n, ... , 1\} but it's not really changing each n deffer from the other and has it's own scope
if you run this code it'll only go so far (around  n = 9613 for this example) until you get a StackOverflowError... not good.

\subsection{Reassignment}
here's a subtle difference, well
hidden behind the overloaded use of the symbol
‘=’, that really sets the two apart. In the imperative
program, ‘=’ refers to a destructive update, assigning
a new value to the left-hand-side variable,
whereas in the functional program ‘=’ means true
equality, and that both the left-hand-side and the right-hand-side can be used interchangeably. This characteristic
of functional programming (known as referential
transparency or purity) has a profound influence
on the way programs are constructed and reasoned
about.

% TODO: My definition:
% before we move on it's important to do the distinction between mutation and Reassignment. simply put when you mutate data the old version of it would become unusable and simple would cease to exist. On the other hand updating a variable should keep the old version usable. a perfect example of this is a VCS (version control system) it gives you the feeling that you are mutating files while in fact it's storing changes (updates) and updating a "HEAD" value to represent the last change (update)



\subsection{Tail Recursion}
Tail recursion provides stack safety to our solutions and prevent Stackoverflows is when you simply return the value of a function call at the end (tail) of your function, in other words your functions has done it's job and handing over the rest of the work to another function
\begin{lstlisting}
  def factorial (n: BigDecimal):BigDecimal = {
    def helper(n: BigDecimal, Acc: BigDecimal): BigDecimal = {
      if (n <= 1) Acc
      else helper(n - 1, n * Acc)
    }
    helper(n, 1)
  }
\end{lstlisting}
Notice that as n takes the values \{n, n-1, n-2, ...\}
the accumulator also changes \{n, n*(n-1), n*(n-1)*(n-2), ...\}
which is very similar to a for loop accumulator pattern, so similar that sometimes compilers compile it to an actual loop


\subsection{Pure Functions}
Pure Functions serves as a (possibly infinite) lookup table mapping from one type to another since variable $x$ will always be the same $f(x)$ will too. this is what's known as referential transparency.
\lstset{style=cpp}
\begin{lstlisting}
    #include <functional>
    #include <iostream>
    int sum(const int v[], const int& n) {
        std::function<int(int, int)> helper = [&v, &n, &helper](int index, int Acc) {
            if (index >= n) return Acc;
            return helper(index + 1, Acc + v[index]);
        };
        return helper(0, 0);
    }
    int main(int, char**) {
        const int a[] = {1, 2, 3, 4, 5, 6};
        // a[1] = 3; not allowed
        int total = sum(a, 6); // 21
        someFunction(a);
        otherFunction(a);
        if (total == sum(a,6))
            std::cout<<"It should be equal same function same argument?!";
        return 0;
    }
\end{lstlisting}
generally speaking $"someFunction"$ and $"otherFunction"$ could've done al sorts of things with $a$ (changing an element value, adding more elements, removing some elements, delete the pointer entirely ...) but if they were pure functions or like in this case using a some sort of a language guarantee (here it's $const$) $a$ will not be modified and in turns $total == sum(a,6)$ and overall our program would be easier to reason about.

\subsection{Laziness}
sometimes referred to as call-by-need, it's the notion that "if data won't change and functions won't neither so would results" meaning that we wouldn't perfome any operations unless they're absolutely necessary 
\lstset{style=scala}
\begin{lstlisting}
    def from(n: Int): LazyList[Int] = 
        n #:: from(n+1) 
        // here from(n+1) won't be evaluated
    def sieve (s: LazyList[Int]): LazyList[Int] = 
        s.head #:: sieve(s.tail.filter(_ % s.head != 0)) 
    // nor would sieve(s.tail.filter(_ % s.head != 0))
    val primes = sieve(from(2))

    primes
        .take(10)
        .toList // now it's necessary

\end{lstlisting}
in this example only 10 elements of the infinite "LazyList" are evaluated while the other elements are immutable they are yet to be evaluated and materialized.

Another huge advantage lazy evaluation gives us


% \subsection{logfilesystem}
% sth
% \subsection{Append Logs}
% sth
% \subsection{Copy-On-Write}
% sth
% \subsection{Structural Sharing}
% effecient cow reduce the overhead of threading and parallel processing
% \subsection{MVC}
% sth
% \subsection{elm}
% sth
% \subsection{Cpp example}
% $https://www.youtube.com/watch?v=y_m0ce1rzRI\&t=4081s$
\section*{References}
[1] 

%TODO: talk about referential transparency
%TODO: use these links:
%TODO: http://typeocaml.com/2015/01/02/immutable/
%TODO: https://algs4.cs.princeton.edu/23quicksort/
%TODO: https://coderscat.com/quicksort-history-and-impls/
%TODO: https://www.codingninjas.com/blog/2020/09/26/mutating-non-mutating-algorithms-in-c/
%TODO: https://elm-lang.org/news/blazing-fast-html
%TODO: https://discourse.elm-lang.org/t/can-the-compiler-skip-virtual-dom/6300
%TODO: https://www.ncbi.nlm.nih.gov/pmc/articles/PMC3136027/
%TODO: https://en.wikipedia.org/wiki/Persistent_data_structure#Garbage_collection
%TODO: https://en.wikipedia.org/wiki/Funarg_problem
%TODO: https://blog.sigplan.org/2022/01/13/provably-space-efficient-parallel-functional-programming/
%TODO: https://en.wikipedia.org/wiki/Persistent_data_structure#Garbage_collection
%TODO: StateManegment
% prefered data structures
% dum data classes
% following strong type system
% funcional
% contexts static compiletime
% curring dynamic clojures
% compile optimizations
% lazyness
% transperncy substitiustion model
% GC or static refs


% chanllenging
% modularity
% time travelling debugger
% bref on 
% determenistic concurrency
% reactiveity
% modernday statless deployment
% 	docker or distributed or actor
% 	strong endpoint devices


\section*{Acknowledgements}
The author wishes to thank Blabla and bla

\printbibliography


\end{document}